
\documentclass[a4paper,11pt]{article}
\usepackage[a4paper, margin=8em]{geometry}

% usa i pacchetti per la scrittura in italiano
\usepackage[french,italian]{babel}
\usepackage[T1]{fontenc}
\usepackage[utf8]{inputenc}
\frenchspacing 

% usa i pacchetti per la formattazione matematica
\usepackage{amsmath, amssymb, amsthm, amsfonts}

% usa altri pacchetti
\usepackage{gensymb}
\usepackage{hyperref}
\usepackage{standalone}

% imposta il titolo
\title{Appunti Comunicazioni Numeriche}
\author{Luca Seggiani}
\date{2026}

% imposta lo stile
% usa helvetica
\usepackage[scaled]{helvet}
% usa palatino
\usepackage{palatino}
% usa un font monospazio guardabile
\usepackage{lmodern}

\renewcommand{\rmdefault}{ppl}
\renewcommand{\sfdefault}{phv}
\renewcommand{\ttdefault}{lmtt}

% disponi il titolo
\makeatletter
\renewcommand{\maketitle} {
	\begin{center} 
		\begin{minipage}[t]{.8\textwidth}
			\textsf{\huge\bfseries \@title} 
		\end{minipage}%
		\begin{minipage}[t]{.2\textwidth}
			\raggedleft \vspace{-1.65em}
			\textsf{\small \@author} \vfill
			\textsf{\small \@date}
		\end{minipage}
		\par
	\end{center}

	\thispagestyle{empty}
	\pagestyle{fancy}
}
\makeatother

% disponi teoremi
\usepackage{tcolorbox}
\newtcolorbox[auto counter, number within=section]{theorem}[2][]{%
	colback=blue!10, 
	colframe=blue!40!black, 
	sharp corners=northwest,
	fonttitle=\sffamily\bfseries, 
	title=~\thetcbcounter: #2, 
	#1
}

% disponi definizioni
\newtcolorbox[auto counter, number within=section]{definition}[2][]{%
	colback=red!10,
	colframe=red!40!black,
	sharp corners=northwest,
	fonttitle=\sffamily\bfseries,
	title=~\thetcbcounter: #2,
	#1
}

% U.D.M
\newcommand{\amp}{\ensuremath{\mathrm{A}}}
\newcommand{\volt}{\ensuremath{\mathrm{V}}}
\newcommand{\meter}{\ensuremath{\mathrm{m}}}
\newcommand{\second}{\ensuremath{\mathrm{s}}}
\newcommand{\farad}{\ensuremath{\mathrm{F}}}
\newcommand{\henry}{\ensuremath{\mathrm{H}}}
\newcommand{\siemens}{\ensuremath{\mathrm{S}}}

% circuiti
\usepackage{circuitikz}
\usetikzlibrary{babel}

% disegni
\usepackage{pgfplots}
\pgfplotsset{width=10cm,compat=1.9}

% disponi codice
\usepackage{listings}
\usepackage[table]{xcolor}

\lstdefinestyle{codestyle}{
		backgroundcolor=\color{black!5}, 
		commentstyle=\color{codegreen},
		keywordstyle=\bfseries\color{magenta},
		numberstyle=\sffamily\tiny\color{black!60},
		stringstyle=\color{green!50!black},
		basicstyle=\ttfamily\footnotesize,
		breakatwhitespace=false,         
		breaklines=true,                 
		captionpos=b,                    
		keepspaces=true,                 
		numbers=left,                    
		numbersep=5pt,                  
		showspaces=false,                
		showstringspaces=false,
		showtabs=false,                  
		tabsize=2
}

\lstdefinestyle{shellstyle}{
		backgroundcolor=\color{black!5}, 
		basicstyle=\ttfamily\footnotesize\color{black}, 
		commentstyle=\color{black}, 
		keywordstyle=\color{black},
		numberstyle=\color{black!5},
		stringstyle=\color{black}, 
		showspaces=false,
		showstringspaces=false, 
		showtabs=false, 
		tabsize=2, 
		numbers=none, 
		breaklines=true
}

\lstdefinelanguage{javascript}{
	keywords={typeof, new, true, false, catch, function, return, null, catch, switch, var, if, in, while, do, else, case, break},
	keywordstyle=\color{blue}\bfseries,
	ndkeywords={class, export, boolean, throw, implements, import, this},
	ndkeywordstyle=\color{darkgray}\bfseries,
	identifierstyle=\color{black},
	sensitive=false,
	comment=[l]{//},
	morecomment=[s]{/*}{*/},
	commentstyle=\color{purple}\ttfamily,
	stringstyle=\color{red}\ttfamily,
	morestring=[b]',
	morestring=[b]"
}

% disponi sezioni
\usepackage{titlesec}

\titleformat{\section}
	{\sffamily\Large\bfseries} 
	{\thesection}{1em}{} 
\titleformat{\subsection}
	{\sffamily\large\bfseries}   
	{\thesubsection}{1em}{} 
\titleformat{\subsubsection}
	{\sffamily\normalsize\bfseries} 
	{\thesubsubsection}{1em}{}

% disponi alberi
\usepackage{forest}

\forestset{
	rectstyle/.style={
		for tree={rectangle,draw,font=\large\sffamily}
	},
	roundstyle/.style={
		for tree={circle,draw,font=\large}
	}
}

% disponi algoritmi
\usepackage{algorithm}
\usepackage{algorithmic}
\makeatletter
\renewcommand{\ALG@name}{Algoritmo}
\makeatother

% disponi numeri di pagina
\usepackage{fancyhdr}
\fancyhf{} 
\fancyfoot[L]{\sffamily{\thepage}}

\makeatletter
\fancyhead[L]{\raisebox{1ex}[0pt][0pt]{\sffamily{\@title \ \@date}}} 
\fancyhead[R]{\raisebox{1ex}[0pt][0pt]{\sffamily{\@author}}}
\makeatother

\begin{document}
% sezione (data)
\section{Lezione del 16-04-26}

% stili pagina
\thispagestyle{empty}
\pagestyle{fancy}

% testo
\subsubsection{Causalità nei SLS}
Definiamo la \textbf{causalità} negi SLS come la proprietà:
\begin{definition}{Causalità}
Un SLS si dice causale se la sua risposta in dominio tempo $h(t)$ rispetta: 
$$
h(t) = h(t) \cdot u(t)
$$
\end{definition}
dove $u(t)$ è il solito gradino di Heaviside. Questa affermazione equivale a dire che il SLS è \textit{causale}, cioè al tempo $t$ la sua uscita non si dipende da informazioni presenti al tempo $t' > t$, o nel futuro. 

\subsubsection{BIBO stabilità nei SLS}
Una seconda proprietà che ci interessa dei SLS è la \textbf{BIBO stabilità}:
\begin{definition}{BIBO stabilità}
Un SLS si dice stabile BIBO se ad ogni ingresso limitato corrisponde un’uscita limitata
\end{definition}

Si può dimostrare che la BIBO stabilità di un sistema SLS equivale a dire:
$$
\int_{-\infty}^{\infty} |h(t)| \, dt \in \mathbb{R} < \infty
$$

\subsection{Campionamento e ricostruzione dei segnali}
La teoria che abbiamo visto finora riguardo alla trasformata di Fourier e i sistemi lineari tempo invarianti (LTI) serve a dimostrare importanti risultati per quanto riguarda il \textit{campionamento} dei segnali in dominio tempo. 

Avevamo introdotto in 3.1.2 l'idea di campionamento di un segnale $x(t)$, ad esempio ottenuto da un microfono, come un metodo che consisteva in "registrare" il valore $x(t)$, per ogni naturale $n$, al tempo $nT$. Si otteneva quindi la funzione:
$$
x[n] = x(nT)
$$

Se $T$ è il \textit{periodo} di campionamento, allora avremo una certa frequenza di campionamento $f_c$ inversa:
$$
f_c = \frac{1}{T}
$$

Nel mondo reale i dispositivi che operano il campionamento sono detti \textbf{ADC}, da \textit{Analog to Digital Converter}. Notiamo che la conversione prevede sia il campionamento del segnale ad ogni istante $nT$, che la quantizzazione dei campioni ottenuti come numeri binari su un certo numero di cifre.

Se il teorema del campionamento che dimostreremo ci darà un limite inferiore per la frequenza di campionamento $f_c$ che \textit{assicurerà}, dal punto di vista teorico, di ricostruire il segnale esattamente, sull'errore di quantizazzione non possiamo fare nulla: questo è irreversibile. 

\begin{enumerate}
	\item 
La prima domanda fondamentale che ci poniamo è quindi la seguente: con quale frequenza dobbiamo campionare un dato segnale $x(t)$ affinché non si perda informazione dal segnale originale, una volta ricostruito?
	\item
La seconda domanda fondamentale che ci poniamo è qual'è l'\textit{interpolatore} da usare per ricostruire il segnale a partire dai campioni ottenuti. In altre parole, ci chiediamo se esiste un filtro con risposta impulsiva $p(t)$ tale per cui:
$$
\hat{x}(t) = \sum_n x(nT) \, p(t - nT) = x(t)
$$

Nel mondo reale gli interpolatori dal dominio digitale a quello analogico vengono detti convertitori \textbf{DAC}, da (\textit{Digital to Analog Converter}).
\end{enumerate}

\subsubsection{Trasformata di Fourier discreta}
In 7.1 avevamo visto la trasformata di Fourier discreta, definita con un tempo di campionamento $T$. In particolare:
\begin{definition}{Trasformata di Fourier discreta}
Data una funzione $x(t)$ campionata $x(nT)$, la sua trasformata di Fourier discreta $\overline{X}(f)$ è:
$$
\overline{X}(f) = \sum_n x(nT) e^{-j2 \pi f n T}
$$
\end{definition}

Anche per questa vale la un \textit{antitrasformata} che ci riporta al domino tempo:
\begin{definition}{Anti di Fourier discreta}
Data una trasformata discreta $\overline{X}(f)$, la sua antitrasformata di Fourier discreta $x(t)$ è: 
$$
x(nT) = T \int_{-\frac{1}{2T}}^{\frac{1}{2T}} \overline{X}(f) e^{-j 2 \pi f n T} \, df 
$$
\end{definition}

Si può dimostrare che:
$$
\overline{X}(f) = \frac{1}{T} \sum_k X \left( f - \frac{k}{T} \right) = f_c \sum_k X(f - k f_c)
$$
dove la $X()$ rappresenta la trasformata continua. Questo significa che la trasformata continua viene \textit{periodicizzata}, e lo spettro si replica (\textit{aliasing}) ai multipli della frequenza di campionamento.

\subsubsection{Condizione di Nyquist}
Per il campionamento di segnali a banda limitata $B$, vale la condizione di aliasing per la rimozione dell'aliasing (e quindi la preservazione del segnale anlogico originale). Qusta si riassume a:
$$
f_c \geq f_n = 2B
$$

In questo, la frequenza minima di campionamento:
$$
f_n = 2B
$$
viene detta \textit{frequenza di Nyquist}.

Per questo motivo quando campioniamo segnali reali (a tempo finito e banda non limitata) implementiamo un \textit{filtro di anti-alias}, inteso come un filtro ideale o brickwall di banda $B'$ scelta appositamente.
Ricordiamo ad esempio che l'orecchio umano è capace di percepire frequenze dai $20-40 \, \mathrm{Hz}$ circa fino ai $20 \, \mathrm{KHz}$, per cui $B' \approx 20 \, \mathrm{KHz}$ è una scelta tipica nell'audio.
Questo filtro rimuoverà quindi l'informazione oltre la frequenza $B'$, portando ad un campionamento valido alla frequenza:
$$
f_c \geq f_n = 2B'
$$
valido, che non introduce aliasing in fase di conversione ADC o DAC. Questo è confermato ad esempio dal fatto che i CD campionano i segnali audio a $44100 \, \mathrm{Hz} \approx 2 \cdot 20 \, \mathrm{KHz}$. 

Questo può essere un buon momento anche per ricordare che la risposta in dominio tempo di un filtro brickwall ideale (in frequenza una $\mathrm{rect}$) equivale alla funzione $\mathrm{sinc}$: questo ci anticipa il fatto che l'interpolatore ideale sarà l'\textit{interpolatore cardinale}, basato sulla $\mathrm{sinc}$ stessa. Avevamo introdotto questo concetto come corollario del teorema 7.4 di dualità, cioè appena avevamo scoperto che la $\mathrm{rect}$ trasformava alla $\mathrm{sinc}$ (e quindi le $\mathrm{sinc}$ fanno da funzioni base dello spazio campionato da $0$ a $B$). 

\end{document}
